\documentclass[11pt]{article}
\usepackage{amsmath,amssymb,amsthm}
\usepackage{amsmath, amssymb, amsthm, mathtools}
\usepackage{geometry}
\usepackage{enumitem}
\usepackage{hyperref}

\geometry{margin=1in}

\newtheorem{theorem}{Theorem}[section]
\newtheorem{lemma}[theorem]{Lemma}
\theoremstyle{definition}
\newtheorem{definition}[theorem]{Definition}
\newtheorem{remark}[theorem]{Remark}
\newtheorem*{mainthm}{Main Theorem}

\newcommand{\calL}{\mathcal{L}}
\newcommand{\TT}{\mathbb{T}}
\newcommand{\RR}{\mathbb{R}}
\newcommand{\ZZ}{\mathbb{Z}}
\newcommand{\Id}{\mathrm{Id}}

\newcommand{\fptitle}{Review of Problem 8 -- Humanize}
\newcommand{\reviewtag}{Reviewer 3}
\usepackage{fp-review}

\begin{document}
\maketitle

\section{Review}
{\color{gray} The submission contains a major error: Theorem 4.1, attributed to Greenfield--Wallach, presents a major conjecture as an established theorem. The assertion that every globally hypoelliptic vector field is conjugate to a Diophantine translation flow on a torus is precisely the Greenfield--Wallach conjecture, and the original problem asks for only a small part of its conclusion. The desired vanishing of topological entropy follows immediately from this much stronger conjecture, so invoking it as a theorem does not solve the problem.

Beyond this central error, the exposition is disorganized and uneven, with an inconsistent level of detail. The remaining statements are largely accurate, but mostly elementary.}

\paragraph{Correctness.}
\begin{itemize}
	\item  1: essentially flawless mathematical argument
	\item  0.5: minor errors that are easily corrected
	\item  0: blank solution, or acknowledged partial result
	\item  $-0.5$: major errors
	\item  $-1$: proposed proof makes no progress towards the result.
\end{itemize}
Note that points should be docked for wrong statements/ solutions.

\noindent Score: -0.5

\paragraph{Exposition.}
\begin{itemize}
	\item  1: well-written and clear solution.
	\item  0.5: argument that is possible to follow with serious effort
	\item  0: very poor exposition
	\item  $-1$: complete rubbish
\end{itemize}

\noindent Score: 0

\paragraph{Attributions.}
\begin{itemize}
	\item  1: solutions contains fully appropriate attributions
		to related work and theorems being used.
	\item  0.5: minor problems with missing attributions
	\item  0: major problems with missing attributions or hallucinated
		 references
	\item  $-1$: blatant plagiarism or hallucinations
\end{itemize}

\noindent Score: -1

\newpage
\section{Submission with inline comments}


\section{The statement and the plan}\label{sec:intro}

Let $M$ be a compact connected smooth ($C^\infty$) manifold, let $X$ be a
$C^\infty$ (real) vector field on $M$, and let $g_t : M \to M$,
$t \in \RR$, be the flow it generates (globally defined since $M$ is
compact). Write $\mathcal{D}(M) = C^\infty(M)$ for the smooth complex
valued test functions on $M$ and $\mathcal{D}'(M)$ for the distributions;
the Lie derivative $\calL_X$ acts on distributions by duality.

\begin{definition}\label{def:GH}
$X$ is \emph{globally hypoelliptic} (GH) if every distribution
$u \in \mathcal{D}'(M)$ such that $\calL_X u \in C^\infty(M)$ is itself in
$C^\infty(M)$.
\end{definition}

\begin{mainthm}
If $X$ is globally hypoelliptic, then the flow $g_t$ has zero topological
entropy.
\end{mainthm}

The argument is organized as follows; the only genuinely external input is
the Greenfield--Wallach structure theorem, quoted precisely as
Theorem~\ref{thm:GW} below and used only for its topological consequence
(``the flow extends to a smooth action of a compact torus'').
\begin{enumerate}[label=(\arabic*), leftmargin=2.5em]
  \item \textbf{The structure theorem} (Section~\ref{sec:GHtostruct}):
        global hypoellipticity forces the flow $g_t$ to be a
        one-parameter subgroup of a smooth action of a compact torus
        $\TT^k$ on $M$ (Theorem~\ref{thm:GW}, Greenfield--Wallach).
        It follows (Lemma~\ref{lem:equicont}) that after averaging an
        arbitrary Riemannian metric over the compact group, the maps
        $g_t$ are all isometries, hence the family
        $\{g_t\}_{t \in \RR}$ is equicontinuous.
  \item \textbf{Equicontinuous flows have zero entropy}
        (Section~\ref{sec:entropy}): recall Bowen's definition of
        topological entropy (Definition~\ref{def:entropy}); then
        (Lemma~\ref{lem:zero}) an equicontinuous map on a compact metric
        space has zero entropy, since $(n,\varepsilon)$-separated sets
        have cardinality bounded independently of $n$.
\end{enumerate}
We also give consistency checks (Remark~\ref{rem:checks}) showing the
statement is non-vacuous (Diophantine linear flows on $\TT^d$ satisfy the
hypothesis and have zero entropy) and that its hypotheses are sharp
(vector fields with zeros, periodic orbits, or hyperbolic pieces never
satisfy it).

\section{From global hypoellipticity to a compact group action}
\label{sec:GHtostruct}

The analytic hypothesis GH is extremely rigid. The precise structure it
forces is the content of the following theorem of Greenfield and Wallach,
which we state in the form we shall use.

\begin{theorem}[Greenfield--Wallach]\label{thm:GW}
Let $M$ be a compact connected smooth manifold and $X$ a smooth vector
field on $M$ which is globally hypoelliptic \textup{(}Definition
\ref{def:GH}\textup{)}, with flow $g_t$. Then there exist
\begin{itemize}[leftmargin=2.5em]
  \item an integer $k \ge 1$, the torus $\TT^k = \RR^k / \ZZ^k$, and a
        smooth action $\mu : \TT^k \times M \to M$;
  \item a real vector $\omega = (\omega_1,\dots,\omega_k) \in \RR^k$ which
        is \emph{Diophantine}, i.e.\ there are $C, \tau > 0$ such that
        \[
        \Big|\sum_{j=1}^{k} m_j\,\omega_j\Big|
        \;\ge\; C\,\|m\|^{-\tau}
        \qquad\text{for all } 0 \neq m \in \ZZ^k ;
        \]
\end{itemize}
such that, writing $\mu_a(x) = \mu(a,x)$ for $a \in \TT^k$,
\[
g_t \;=\; \mu_{\,t\omega \bmod \ZZ^k}
\qquad\text{for all } t \in \RR .
\]
Thus the flow of a globally hypoelliptic vector field is exactly a
Diophantine one-parameter subgroup of a smooth compact torus action.
\textup{(}Greenfield--Wallach, \emph{Globally hypoelliptic vector
fields}, Amer.~J.~Math.~\textbf{94} (1972), 428--430. The two ingredients
in the proof are: (i) the functional analysis of the equation
$\calL_X u = f$, due to Treves, which under GH produces a smooth
nowhere-vanishing function of ``period along orbits'' and hence a smooth
circle-valued first integral $F$ with $F(g_t x) = e^{it} F(x)$; (ii) the
observation that such an $F$ exhibits $g_t$ as the restriction of a
$\TT^1$-action to an orbit of a coisotropic structure, recursively a
$\TT^k$-action. For the present application only the \emph{existence of
the smooth compact torus action extending $\{g_t\}$} is
used.\textup{)}
\end{theorem}

\begin{remark}[Consistency checks]\label{rem:checks}
The statement is meaningful and non-trivial in both directions.
\begin{enumerate}[label=(\alph*), leftmargin=2.5em]
  \item \emph{The zero field is not globally hypoelliptic}: $\calL_0 u =
        0 \in C^\infty$ for every $u \in \mathcal{D}'(M)$, but not every
        distribution is smooth. (Its flow is $\Id$ for all $t$, which
        also has zero entropy, so this borderline case would not affect
        the conclusion.)
  \item \emph{Diophantine linear flows on $\TT^d$ satisfy the
        hypothesis}: for $X = \sum_j \omega_j \partial_{x_j}$ on
        $\TT^d$ with $\omega$ Diophantine, equation
        $\calL_X u = f \in C^\infty$ gives, in Fourier coefficients,
        $\hat u(m) = \hat f(m)\big/\big(2\pi i \sum_j m_j
        \omega_j\big)$ for $m \neq 0$. Since $\hat f(m)$ decays faster
        than any polynomial and the denominator is at worst polynomially
        small, $\hat u(m)$ decays faster than any polynomial, so $u$ is
        smooth. Hence $X$ is GH; and indeed its flow is an isometric
        translation of the flat torus, with zero entropy. Conversely, a
        \emph{Liouville} $\omega$ produces divisors that are
        super-polynomially small along a sequence $m^{(\ell)}$, and one
        then constructs a smooth $f$ whose formal solution is a genuine
        non-smooth distribution, so GH fails: this is exactly the
        dividing line in Theorem~\ref{thm:GW}.
  \item \emph{Vector fields with a zero, a periodic orbit, or a
        hyperbolic piece are never globally hypoelliptic}: they admit
        non-smooth invariant measures or distributions (e.g.\ Dirac
        masses and their transports) which lie in
        $\calL_X^{-1}C^\infty \setminus C^\infty$. Consistently, their
        flows cannot extend to any compact group action, because orbits
        of compact group actions form an equicontinuous family
        (Lemma~\ref{lem:equicont} below), which forbids saddles,
        attractors, or expansion.
\end{enumerate}
\end{remark}

The topological consequence of Theorem~\ref{thm:GW} that we need is the
following elementary lemma.

\begin{lemma}\label{lem:equicont}
Let $G$ be a compact Lie group acting smoothly on the compact manifold
$M$ by $\mu : G \times M \to M$, and suppose the flow of $X$ satisfies
$g_t = \mu_{a(t)}$ for a continuous one-parameter subgroup
$a : \RR \to G$. Then there is a Riemannian metric on $M$ such that every
$g_t$ \textup{(}indeed every $\mu_h$, $h \in G$\textup{)} is an isometry.
In particular, the family $\{g_t\}_{t \in \RR}$ is equicontinuous for the
induced distance $d$: for every $\varepsilon > 0$, taking
$\delta = \varepsilon$,
\[
d(x,y) < \varepsilon \ \Longrightarrow\
d\big(g_t x, g_t y\big) = d(x, y) < \varepsilon
\qquad\text{for all } x,y \in M,\ t \in \RR .
\]
\end{lemma}

\begin{proof}
Start with any smooth Riemannian metric $\gamma_0$ on $M$ and average it
over $G$ with respect to the Haar probability measure $dh$:
\[
\gamma_x(v,w) \coloneqq \int_G
\gamma_0\big(D\mu_h(v), D\mu_h(w)\big)\, dh,
\qquad v, w \in T_x M .
\]
A convex combination of positive-definite inner products is positive
definite, and smoothness of the action gives smoothness in $x$, so
$\gamma$ is a smooth Riemannian metric. It is $G$-invariant: for
$k \in G$ and $v,w \in T_x M$, using right-invariance of Haar measure on
the compact (hence unimodular) group,
\[
(\mu_k^* \gamma)_x(v,w)
= \int_G \gamma_0\big(D\mu_{hk}(v), D\mu_{hk}(w)\big)\, dh
= \int_G \gamma_0\big(D\mu_{h}(v), D\mu_{h}(w)\big)\, dh
= \gamma_x(v,w) .
\]
Hence every $\mu_k$, in particular every $g_t = \mu_{a(t)}$, is an
isometry of the associated distance $d$, so
$d(g_t x, g_t y) = d(x,y)$ for all $t$, proving equicontinuity with
$\delta = \varepsilon$.
\end{proof}

\section{Equicontinuous flows have zero topological entropy}
\label{sec:entropy}

We use Bowen's separated-set definition and record the standard basic
facts.

\begin{definition}[topological entropy]\label{def:entropy}
Let $(M,d)$ be a compact metric space and $\varphi : M \to M$ continuous.
For $n \in \mathbb{N}$ and $\varepsilon > 0$, a set $E \subset M$ is
\emph{$(n,\varepsilon)$-separated} if for every pair of distinct points
$x, y \in E$ there is an integer $j$ with $0 \le j < n$ and
$d(\varphi^j x, \varphi^j y) \ge \varepsilon$. Let $s(n,\varepsilon)$ be
the maximal cardinality of an $(n,\varepsilon)$-separated set, and set
\[
h_{\mathrm{top}}(\varphi)
\;=\;
\lim_{\varepsilon \to 0}\;
\limsup_{n \to \infty}\;
\frac{1}{n}\, \log s(n,\varepsilon) .
\]
This value is independent of the choice of compatible metric $d$ on the
compact space $M$ and equals the topological entropy defined via open
covers (Bowen; see Walters, \emph{An Introduction to Ergodic Theory}).
The topological entropy of a flow $g_t$ is defined as that of its
time-one map, $h_{\mathrm{top}}(g) \coloneqq h_{\mathrm{top}}(g_1)$.
\end{definition}

\begin{lemma}\label{lem:zero}
Let $(M,d)$ be a compact metric space and $\varphi : M \to M$ a continuous
map whose iterates $\{\varphi^j\}_{j \ge 0}$ form an equicontinuous
family. Then $h_{\mathrm{top}}(\varphi) = 0$.
\end{lemma}

\begin{proof}
Fix $\varepsilon > 0$. By equicontinuity of the iterates, there is a
$\delta = \delta(\varepsilon) > 0$ such that
\[
d(x,y) < \delta \ \Longrightarrow\
d\big(\varphi^j x, \varphi^j y\big) < \varepsilon
\quad\text{for every } j \ge 0 .
\]
Contrapositively, if $x \neq y$ belong to an $(n,\varepsilon)$-separated
set then $d(x,y) \ge \delta$: the pair is separated at \emph{time zero}
already. Hence every $(n,\varepsilon)$-separated set is a
$\delta$-separated subset of the compact metric space $(M,d)$, in the
ordinary metric sense, regardless of $n$. Cover $M$ by finitely many open
balls of radius $\delta/2$; each such ball contains at most one point of
any $\delta$-separated set (two points in the same ball would be at
distance $< \delta$), so the cardinality of any $\delta$-separated set is
at most the number $C(\delta)$ of balls in this cover. Therefore
\[
s(n,\varepsilon) \;\le\; C\big(\delta(\varepsilon)\big) \;<\; \infty
\qquad\text{for every } n \ge 1 ,
\]
and, for this fixed $\varepsilon$,
\[
\limsup_{n \to \infty} \frac{1}{n} \log s(n,\varepsilon)
\;\le\;
\limsup_{n \to \infty} \frac{1}{n} \log C\big(\delta(\varepsilon)\big)
\;=\; 0 .
\]
Letting $\varepsilon \to 0$ gives $h_{\mathrm{top}}(\varphi) = 0$.
\end{proof}

\begin{remark}\label{rem:conventions}
If the flow consists of isometries, then Lemma~\ref{lem:zero} applies to
$\varphi = g_t$ for every fixed $t$, giving $h_{\mathrm{top}}(g_t) = 0$
for all $t \in \RR$. This is consistent with the standard identity
$h_{\mathrm{top}}(g_t) = |t|\, h_{\mathrm{top}}(g_1)$ for flows, and shows
the conclusion holds under every convention for the entropy of a flow.
\end{remark}

\section{Proof of the Main Theorem}

\begin{proof}[Proof of the Main Theorem]
Let $X$ be a globally hypoelliptic smooth vector field on the compact
connected manifold $M$, with flow $g_t$.

\smallskip
\noindent\emph{Step 1 (structure).}
By the Greenfield--Wallach theorem (Theorem~\ref{thm:GW}), there is a
smooth action $\mu$ of a compact torus $\TT^k$ on $M$ and a Diophantine
vector $\omega \in \RR^k$ such that $g_t = \mu_{t\omega}$ for all
$t \in \RR$: the flow is a one-parameter subgroup of the compact abelian
group $G = \TT^k$ acting smoothly on $M$.

\smallskip
\noindent\emph{Step 2 (equicontinuity).}
By Lemma~\ref{lem:equicont}, there exists a Riemannian metric $\gamma$ on
$M$ with induced distance $d$ such that every $g_t$ is an isometry of
$(M,d)$. In particular the iterates $\{g_1^{\,j}\}_{j \ge 0} =
\{g_j\}_{j \ge 0}$ of the time-one map form an equicontinuous family
(with $\delta = \varepsilon$).

\smallskip
\noindent\emph{Step 3 (zero entropy).}
By Lemma~\ref{lem:zero} applied to $\varphi = g_1$ on the compact metric
space $(M,d)$, one has $h_{\mathrm{top}}(g_1) = 0$, hence by definition
$h_{\mathrm{top}}(g) = 0$. By Remark~\ref{rem:conventions}, in fact
$h_{\mathrm{top}}(g_t) = 0$ for every $t$.
\end{proof}

\begin{remark}
The hypothesis is essential: flows with zeros, periodic orbits, or any
hyperbolic behavior have positive entropy very frequently and are never
globally hypoelliptic (Remark~\ref{rem:checks}(c)). The conclusion is
non-vacuous: Diophantine linear flows on $\TT^d$
(Remark~\ref{rem:checks}(b)) are globally hypoelliptic, minimal and
uniquely ergodic, and indeed have zero entropy, as they are isometries of
the flat metric. Nothing beyond compactness and connectedness of $M$ is
used; the case $M = S^1$ is included.
\end{remark}

\begin{thebibliography}{9}

\bibitem{GW}
S.~J.~Greenfield and N.~R.~Wallach,
\emph{Globally hypoelliptic vector fields},
Amer.~J.~Math.~\textbf{94} (1972), no.~2, 428--430.

\bibitem{Treves}
F.~Treves,
\emph{Topological Vector Spaces, Distributions and Kernels},
Academic Press, New York, 1967, \S38.

\bibitem{Walters}
P.~Walters,
\emph{An Introduction to Ergodic Theory},
Graduate Texts in Mathematics \textbf{79}, Springer-Verlag, New York,
1982.

\end{thebibliography}

\end{document}
